\chapter{Extension spectral sequences}
\label{chap:extension-ss}

This chapter develops the general extension-spectral-sequence construction and its internal formal properties.  Any cited factorization theorem is kept as a separate external result.

\section{Construction and elementwise notation}

% SOURCE: aimpaper/main.tex:253-271, Definition def:ess.
\begin{definition}[Filtered two-term complex]
  \label{def:filtered-two-term-complex}
  \uses{def:admissible-spectrum,def:adams-filtration,def:mathlib-chain-complex}
  \notready
  For a stable map $f:X\to Y$, let $C(f)$ be the two-term chain complex
  \[
    0\longrightarrow\pi_*X\xrightarrow{\pi_*f}\pi_*Y
      \longrightarrow0,
  \]
  with both nonzero terms filtered by Adams filtration.  Its homology is
  $\ker(\pi_*f)\oplus\operatorname{coker}(\pi_*f)$.  The filtered-complex
  structure records the fact that $f$ does not lower Adams filtration.
\end{definition}

% SOURCE: aimpaper/main.tex:264-271, Definition def:ess.
\begin{definition}[$f$-extension spectral sequence]
  \label{def:f-extension-ss}
  \uses{def:filtered-two-term-complex,def:core-filtered-chain-complex,def:ess-complex-shape,thm:filtered-complex-to-spectral-sequence,def:ss-einfty-presentation,def:ss-convergence-detection-witness}
  \notready
  The $f$-extension spectral sequence, or $f$-ESS, is the spectral sequence of
  the Adams-filtered two-term complex.  Given convergence-and-detection
  witnesses for the source and target Adams sequences, its initial page and
  abutment comparison are
  \[
    {}^fE_0^{s,t}\cong E_\infty^{s,t}(X)\oplus
      E_\infty^{s,t}(Y)
    \Longrightarrow
    \ker(\pi_*f)\oplus\operatorname{coker}(\pi_*f),
  \]
  and its differential has degree
  $d_n^f:{}^fE_n^{s,t}\to{}^fE_n^{s+n,t+n}$.  The $d_0$ map is induced by
  $f$ on associated graded homotopy groups.  This comparison is a
  project-specific witness, not a consequence of a general filtered-complex
  convergence theorem.
\end{definition}

% SOURCE: aimpaper/main.tex:253-271; project-specific comparison required by
% the Section 2 ESS arguments.  The source and target convergence witnesses
% are explicit inputs rather than a generic filtered-complex theorem.
\begin{theorem}[$f$-ESS initial page and abutment comparison]
  \label{thm:fess-e0-abutment-comparison}
  \uses{def:ss-convergence-detection-witness,def:filtered-two-term-complex}
  \notready
  Let $f:X\to Y$ be a filtration-preserving stable map, and give the classical
  Adams sequences of $X$ and $Y$ convergence-and-detection witnesses.  The
  filtered two-term complex defining the $f$-ESS has
  \[
    {}^fE_0^{s,t}\cong E_\infty^{s,t}(X)\oplus E_\infty^{s,t}(Y),
  \]
  and its abutment is naturally
  \[
    H_*C(f)\cong\ker(\pi_*f)\oplus\operatorname{coker}(\pi_*f).
  \]
  The comparison is compatible with the induced $d_0$ map and with the
  representative ambiguity in both convergence witnesses.  This is the only
  abutment comparison required by the Section 2 ESS proofs; it does not claim
  a general strong-convergence theorem for arbitrary filtered complexes.
\end{theorem}

\begin{proof}
  Apply the two specified witnesses to identify the associated graded terms of
  the two filtered homotopy groups with their algebraic $E_\infty$ pages.
  Compute the homology of the two-term complex directly as kernel plus
  cokernel, and check that the induced map on associated graded objects is the
  $d_0$ differential.  Naturality of the witness comparisons gives the stated
  compatibility.
\end{proof}

\section{Naturality, composition, and exactness}

% SOURCE: aimpaper/main.tex:461-501, Theorem thm:4114f70c.
\begin{theorem}[Commutative-square propagation]
  \label{thm:fess-commutative-square}
  \uses{prop:no-crossing-iff-uniform-detection}
  \notready
  Let
  \[
  \begin{array}{ccc}
    X&\xrightarrow{f}&Y\\
    \mathord{p}\downarrow&&\downarrow\mathord{q}\\
    Z&\xrightarrow{g}&W
  \end{array}
  \qquad(qf=gp)
  \]
  be homotopy commutative.  Fix $m,n,l\ge0$ and
  $0<k\le m+l-n$, and classes
  \[
  \begin{aligned}
    x&\in E_\infty^{s,t}(X),&
    y&\in E_\infty^{s+n,t+n}(Y),\\
    z&\in E_\infty^{s+m,t+m}(Z),&
    w&\in E_\infty^{s+m+l,t+m+l}(W).
  \end{aligned}
  \]
  Suppose $d_n^f(x)=y$, $d_m^p(x)=z$, at least one of these two extensions
  has no crossing, and $d_l^g(z)=w$ has no crossing in filtrations at least
  $s+n+k$.  If the zero extension $d_{k-1}^q(y)=0$ has no crossing, then
  \[
    d_{m+l-n}^q(y)=w.
  \]
\end{theorem}

\begin{proof}
  Use no-crossing on one of the two source extensions to choose a single
  representative $[x]$ realizing both $y$ and $z$.  The zero $q$-extension
  forces $qf[x]$ into filtration at least $s+n+k$.  Commutativity identifies
  it with $gp[x]$, and no-crossing for the $g$-extension identifies its leading
  term with $w$.  Apply the representative characterization once more to get
  the asserted $q$-extension.
\end{proof}

% SOURCE: aimpaper/main.tex:555-574, Corollary cor:0012nik.
\begin{corollary}[Unrestricted square propagation]
  \label{cor:fess-square-naturality}
  \uses{thm:fess-commutative-square}
  \notready
  In the same square, if $d_n^f(x)=y$, $d_m^p(x)=z$, one of these has no
  crossing, and $d_l^g(z)=w$ has no crossing, then
  $d_{m+l-n}^q(y)=w$ whenever the displayed length is nonnegative.
\end{corollary}

\begin{proof}
  Apply Theorem~\ref{thm:fess-commutative-square} at the first filtration where
  a nonzero $q$-extension could occur.  Complete no-crossing for the
  $g$-extension supplies every lower-bound hypothesis.
\end{proof}

% SOURCE: aimpaper/main.tex:576-595, Corollary cor:166dc180.
\begin{corollary}[Triangle-shaped shift]
  \label{cor:fess-triangle-shift}
  \uses{cor:fess-square-naturality}
  \notready
  If $p=qf$, $d_n^f(x)=y$, and $d_m^p(x)=z$, and at least one of these
  extensions has no crossing, then $d_{m-n}^q(y)=z$ whenever $m\ge n$.
\end{corollary}

\begin{proof}
  Use the unrestricted square theorem with the bottom map the identity and
  its length-zero extension.
\end{proof}

% SOURCE: aimpaper/main.tex:597-615, Corollary cor:290d35ce.
% VERIFIED-IN: main.tex:604 has z in E_infty(Y); the diagram requires E_infty(Z).
\begin{corollary}[Composition of successive extensions]
  \label{cor:fess-composition}
  \uses{cor:fess-square-naturality}
  \notready
  If $q=gp$, $d_m^p(x)=z$ with
  $z\in E_\infty^{s+m,t+m}(Z)$, and $d_l^g(z)=w$ has no crossing, then
  $d_{m+l}^q(x)=w$.  The target spectrum of $z$ is $Z$; the occurrence of
  $E_\infty(Y)$ in \texttt{aimpaper/main.tex:604} is a corrected source typo.
\end{corollary}

\begin{proof}
  Apply square naturality with the top map the identity and its length-zero
  extension.
\end{proof}

% SOURCE: aimpaper/main.tex:617-630, Corollary cor:e7b20ae2.
\begin{corollary}[Stable-page maps between extension spectral sequences]
  \label{cor:fess-map-of-stable-pages}
  \uses{cor:fess-square-naturality,def:reindexed-ss-morphism}
  \notready
  For a commutative square $qf=gp$, suppose the $p$- and $q$-ESS have no
  differential before page $r$.  Then $(d_r^p,d_r^q)$ induces compatible page
  maps from the $f$-ESS to the $g$-ESS, commuting with every later extension
  differential with the required reindexing.
\end{corollary}

\begin{proof}
  Initial-page stability gives no-crossing for $d_r^p$ and $d_r^q$.  Apply
  square naturality first to $d_0^f$ and $d_0^g$, then induct over ESS pages.
  At each stage the commuting differential square descends to homology and
  constructs the next page map.
\end{proof}

% SOURCE: aimpaper/main.tex:632-646, Corollary cor:aed3d1a4.
\begin{corollary}[A zero composite makes the target permanent]
  \label{cor:fess-zero-composite-permanent}
  \uses{prop:fextension-iff-detected,cor:fess-square-naturality}
  \notready
  If $gf=0$ and $d_n^f(x)=y$, then $y$ is a permanent cycle in the $g$-ESS:
  $d_m^g(y)=0$ for every $m\ge0$.
\end{corollary}

\begin{proof}
  Choose $[x]$ with $f[x]$ detected by $y$.  Its image under $g$ is zero, so
  no nonzero leading term can occur in the $g$-ESS.  Equivalently, apply square
  naturality to the square with zero bottom and left maps.
\end{proof}

% SOURCE: aimpaper/main.tex:648-664, Proposition prop:cfe810af.
% VERIFIED-IN: only exactness at pi_*Y is used; the zeros printed at both ends are not assumptions.
\begin{proposition}[Exactness at the middle makes every permanent cycle a boundary]
  \label{prop:fess-exact-middle-boundary}
  \uses{def:f-extension-ss,prop:fextension-iff-detected,cor:fess-zero-composite-permanent}
  \notready
  If $X\xrightarrow fY\xrightarrow gZ$ induces a sequence exact at
  $\pi_*Y$, then every permanent $g$-ESS cycle in $E_\infty(Y)$ is hit by an
  $f$-ESS differential.  No injectivity of $f$ or surjectivity of $g$ is
  assumed.
\end{proposition}

\begin{proof}
  A permanent $g$-ESS cycle has a detected representative $[y]$ with
  $g[y]=0$.  Exactness at the middle supplies $[x]$ with $f[x]=[y]$.
  Choose the leading Adams class $x$ detecting $[x]$ and apply the
  representative characterization of $f$-extensions.  This avoids the
  stronger short-exact sequence accidentally displayed in the paper's first
  proof.
\end{proof}

\section{Representatives, essentiality, and crossings}

% SOURCE: aimpaper/main.tex:277-295, Notation nota:6fb333a2.
\begin{definition}[$f$-ESS cycles and boundaries]
  \label{def:fess-cycles-boundaries}
  \uses{def:f-extension-ss}
  \notready
  Define ${}^fZ_n^{s,t}(X)\subseteq E_\infty^{s,t}(X)$ to consist of classes
  on which $d_0^f,\ldots,d_n^f$ vanish, and let
  ${}^fB_n^{s,t}(Y)$ be the subgroup generated by their images.  Put
  ${}^fZ_{-1}=E_\infty(X)$ and ${}^fB_{-1}=0$.  Then
  \[
    {}^fE_n^{s,t}\cong{}^fZ_{n-1}^{s,t}(X)\oplus
      E_\infty^{s,t}(Y)/{}^fB_{n-1}^{s,t}(Y),
  \]
  and $d_n^f(x)=y$ means equality with the coset of $y$ modulo
  ${}^fB_{n-1}$.
\end{definition}

% SOURCE: aimpaper/main.tex:297-304, Definition def:768fba8a and following notation.
\begin{definition}[$f$-extension and essentiality]
  \label{def:f-extension-essential}
  \uses{def:fess-cycles-boundaries}
  \notready
  An $f$-extension of length $n$ from
  $x\in E_\infty^{s,t}(X)$ to
  $y\in E_\infty^{s+n,t+n}(Y)$ is the assertion $d_n^f(x)=y$.
  It is essential exactly when
  $y\notin{}^fB_{n-1}^{s+n,t+n}(Y)$; otherwise it is inessential.  For a
  permanent class $x$, the notation $\{x\}$ denotes the set of homotopy
  classes detected by $x$, and $[x]$ denotes an explicitly chosen member of
  that set.
\end{definition}

% SOURCE: aimpaper/main.tex:307-325, Proposition prop:8154e6f1(1).
\begin{proposition}[Extension iff a representative maps to a representative]
  \label{prop:fextension-iff-detected}
  \uses{def:f-extension-essential,def:adams-detection}
  \notready
  There is an $f$-extension $d_n^f(x)=y$ if and only if some
  $[x]\in\{x\}$ satisfies $f[x]\in\{y\}$, where the target filtration is
  $s+n$ and equality is understood modulo ${}^fB_{n-1}$.
\end{proposition}

\begin{proof}
  In the filtered two-term complex, choose a lift of $x$ in filtration $s$.
  Its image has first nonzero associated-graded term in filtration $s+n$
  exactly when the $n$th ESS differential is the coset of $y$.  Conversely,
  a detected target representative determines such a filtered lift.  Verify
  that changing either lift changes the target only by a shorter ESS boundary.
\end{proof}

% SOURCE: aimpaper/main.tex:307-325, Proposition prop:8154e6f1(2)-(3).
\begin{proposition}[Inessential extensions and competing targets]
  \label{prop:fextension-inessential-iff-higher}
  \uses{prop:fextension-iff-detected,def:fess-cycles-boundaries}
  \notready
  Let $d_n^f(x)=y$.  It is inessential if and only if there are
  $0<a\le n$ and $x'\in E_\infty^{s+a,t+a}(X)$ with an essential extension
  $d_{n-a}^f(x')=y$.  If the same source $x$ has targets $y$ and $y'$ of the
  same filtration, then $y-y'\in{}^fB_{n-1}$; when $y-y'\ne0$, it is the
  target of a shorter essential extension from some higher-filtration source.
  If $y=y'$, only the boundary-membership conclusion is asserted.  If the
  filtration of $y'$ is strictly larger than that of $y$, the extension to
  $y$ is inessential.
\end{proposition}

\begin{proof}
  Expand membership in ${}^fB_{n-1}$ as a finite sum of shorter differential
  images and take the summand of smallest source filtration.  Conversely, any
  shorter essential image belongs to ${}^fB_{n-1}$.  For competing targets,
  subtract the two representative equations and apply the same leading-term
  argument; over $\mathbb F_2$ subtraction is addition, but the statement is
  kept additive.
\end{proof}

% SOURCE: aimpaper/main.tex:344-352, Definition def:98skj23.
\begin{definition}[Crossing and no-crossing range]
  \label{def:fess-crossing}
  \uses{def:f-extension-essential}
  \notready
  A crossing of $d_n^f(x)=y$, with $x$ in filtration $s$, is an extension
  $d_m^f(x')=y'\ne0$ whose source has filtration $s+a$ for $a>0$ and whose
  target filtration lies between a specified lower bound $p$ and
  $\operatorname{AF}(y)=s+n$.  The original extension has no crossing in
  filtrations at least $p$ when no such extension exists; it has no crossing
  without qualification when $p=s+1$.
\end{definition}

% SOURCE: aimpaper/main.tex:403-405, Remark 2.9.
\begin{lemma}[Every crossing contains an essential crossing]
  \label{lem:crossing-has-essential-crossing}
  \uses{def:fess-crossing,prop:fextension-inessential-iff-higher}
  \notready
  If an extension has a crossing in a fixed filtration interval, then it has
  an essential crossing, possibly of smaller length, whose target remains in
  the same interval.
\end{lemma}

\begin{proof}
  If the given crossing is inessential, replace it using
  Proposition~\ref{prop:fextension-inessential-iff-higher}.  The replacement
  has larger source filtration and the same target.  Length strictly
  decreases, so finite descent terminates at an essential crossing without
  leaving the target-filtration interval.
\end{proof}

% SOURCE: aimpaper/main.tex:408-427, Proposition prop:i8r47oe.
\begin{proposition}[No crossing iff uniform detection]
  \label{prop:no-crossing-iff-uniform-detection}
  \uses{def:fess-crossing,prop:fextension-iff-detected,prop:fextension-inessential-iff-higher,lem:crossing-has-essential-crossing}
  \notready
  An extension from $x$ to $y$ has no crossing in filtrations at least $p$ if
  and only if every $[x]\in\{x\}$ satisfying
  $\operatorname{AF}(f[x])\ge p$ has $f[x]\in\{y\}$.  In particular:
  \begin{enumerate}
    \item ordinary no-crossing is equivalent to this conclusion for every
      representative of $x$;
    \item $d_n^f(x)=0$ has no crossing exactly when every representative image
      has filtration strictly greater than $\operatorname{AF}(x)+n$;
    \item if $\operatorname{AF}(f)=n$, every $d_n^f$ has no crossing.
  \end{enumerate}
\end{proposition}

\begin{proof}
  If a representative has a different leading target, compare it with a
  representative producing $y$ and use the competing-target proposition to
  obtain a shorter crossing.  Conversely, add to $[x]$ a higher-filtration
  representative producing the crossing target; its image has a different
  leading term, contradicting uniform detection.  The three special cases
  follow by choosing $p=s+1$, taking the zero coset, and using the degree bound
  for a map of exact filtration $n$.
\end{proof}

% SOURCE: aimpaper/main.tex:430-456, Example 2.11; underlying values cited to IWX.
\begin{proposition}[$\eta$-ESS regression in stem $46$]
  \label{prop:eta-ess-regression}
  \uses{def:f-extension-essential,def:fess-crossing,def:external-evidence}
  \lean{KIP126.Classical.Regression.etaEss}
  \notready
  With explicit IWX computation evidence, the four essential
  $\eta$-extensions in Example 2.11, the inessential extension
  $d_2^\eta(h_0h_3^2h_5)=h_1h_5d_0$, and the stated crossing of
  $d_4^\eta(h_3^2h_5)=Mh_1$ are reproduced.  The regression also checks that
  the length-one extension from $h_5d_0$ has no crossing.
\end{proposition}

\begin{proof}
  Load each typed class, bidegree, and extension from the source-bearing
  computation record.  Evaluate Definitions~\ref{def:f-extension-essential}
  and~\ref{def:fess-crossing}; the shorter length-one extension witnesses both
  inessentiality and the crossing, while the map-filtration criterion proves
  the final no-crossing assertion.
\end{proof}
